\documentclass[UTF8,fontset=fandol,11pt]{ctexart}
\usepackage[a4paper,margin=24mm]{geometry}
\usepackage{amsmath,amssymb,amsthm,graphicx,tabularx,booktabs,hyperref}
\newtheorem{definition}{定义}[section]
\newtheorem{theorem}[definition]{定理}
\newtheorem{example}[definition]{例题}
\newtheorem{exercise}[definition]{练习}
\hypersetup{unicode=true,colorlinks=true,linkcolor=blue}
\title{数列与数列极限：经济学本科一年级微积分讲义计划}
\date{}
\begin{document}
\maketitle
\section{数列与数列极限的定义}\label{sec:S01}
% UNIT S01.U01
\begin{definition}
\label{def:sequence}数列是定义在正整数集 $\mathbb{N}^{+}=\{1,2,3,\ldots\}$ 上、取实数值的函数。若该函数在 $n\in\mathbb{N}^{+}$ 处的值为 $a_n$，则称 $a_n$ 为数列 $\{a_n\}$ 的第 $n$ 项，通常写作
\[
a_1,a_2,\ldots,a_n,\ldots,
\qquad a_n=f(n),\quad n\in\mathbb{N}^{+}.
\]
数列的下标只能取正整数，且项的排列顺序是数列定义的一部分。
\end{definition}
% UNIT S01.U02
\begin{definition}
\label{def:sequence-limit}\label{def:epsilon-N}设 $\{a_n\}$ 是实数数列，$L\in\mathbb{R}$。若对于任意 $\varepsilon>0$，都存在正整数 $N$，使得当 $n\geq N$ 时有
\[
|a_n-L|<\varepsilon,
\]
则称数列 $\{a_n\}$ 收敛于 $L$，记为
\[
\lim_{n\to\infty}a_n=L.
\]
其中，$\varepsilon$ 先任意给定，而 $N$ 可以依赖于 $\varepsilon$，但不能依赖于 $n$。
\end{definition}
% UNIT S01.U03
\begin{definition}
“收敛于给定的 $L$”的否定形式，即“不收敛于给定的 $L$”，是：存在某个 $\varepsilon_0>0$，使得对每一个正整数 $N$，都存在 $n\geq N$，满足
\[
|a_n-L|\geq\varepsilon_0.
\]
因此，收敛于这个给定的 $L$ 要求从某个下标开始的所有项都落在区间 $(L-\varepsilon,L+\varepsilon)$ 内；不收敛于这个给定的 $L$ 则表示无论下标取得多大，总能找到一项落在某个固定误差范围之外。
\end{definition}
% UNIT S01.U04
\begin{theorem}
\label{thm:limit-unique}若数列 $\{a_n\}$ 收敛于 $L$，又收敛于 $M$，其中 $L,M\in\mathbb{R}$，则
\[
L=M.
\]
\end{theorem}
% UNIT S01.U05
\begin{proof}
设 $L\ne M$，取
\[
\varepsilon=\frac{|L-M|}{3}>0.
\]
由 $a_n\to L$，存在正整数 $N_1$，使得当 $n\geq N_1$ 时，$|a_n-L|<\varepsilon$。同理由 $a_n\to M$，存在正整数 $N_2$，使得当 $n\geq N_2$ 时，$|a_n-M|<\varepsilon$。令 $N=\max\{N_1,N_2\}$，则对任意 $n\geq N$，三角不等式给出
\[
|L-M|\leq |L-a_n|+|a_n-M|<2\varepsilon=\frac{2}{3}|L-M|.
\]
这与 $|L-M|>0$ 矛盾。因此 $L=M$。
\end{proof}
% UNIT S01.U06
\begin{example}
验证数列
\[
a_n=\frac{1}{n},\qquad n\in\mathbb{N}^{+},
\]
是否收敛，并判断其极限。
\end{example}
% UNIT S01.U07
\par\noindent\textbf{解：}要证明 $a_n\to0$，需验证对任意 $\varepsilon>0$，存在正整数 $N$，使得 $n\geq N$ 时
\[
\left|\frac{1}{n}-0\right|<\varepsilon.
\]
给定 $\varepsilon>0$，取正整数 $N>\frac{1}{\varepsilon}$。当 $n\geq N$ 时，由 $n>0$ 得
\[
\left|\frac{1}{n}-0\right|=\frac{1}{n}\leq\frac{1}{N}<\varepsilon.
\]
所以按照 $\varepsilon$--$N$ 定义，
\[
\lim_{n\to\infty}\frac{1}{n}=0.
\]\par
% UNIT S01.U08
极限描述的是下标充分大时数列项的接近趋势，而不是某一项必须等于极限。例如，对 $a_n=1/n$，没有正整数 $n$ 使得 $a_n=0$，但当 $n$ 足够大时，$a_n$ 可以小于任意预先给定的正数 $\varepsilon$。因此，有限项的取值与数列的极限需要区分。\par
% UNIT S01.U09
\begin{exercise}
\label{mathagent:exercise:41706b9e371cbf8283a4f1192e660b4426644ffb8f650af52f47cef5491b1ca0}
用 $\varepsilon$--$N$ 定义证明
\[
\lim_{n\to\infty}\frac{2n+1}{n}=2.
\]
要求明确写出给定的 $\varepsilon$、所取的 $N$，以及当 $n\geq N$ 时的误差估计。
\end{exercise}
\section{收敛数列的基本性质与极限运算法则}\label{sec:S02}
% UNIT S02.U01
\begin{definition}
设数列 $\{a_n\}$。若存在常数 $M>0$，使得对所有 $n\in\mathbb{N}^{+}$ 都有 $|a_n|\leq M$，则称数列 $\{a_n\}$ 有界。
\end{definition}
% UNIT S02.U02
\begin{theorem}
\label{thm:convergent-bounded}若数列 $\{a_n\}$ 收敛于实数 $L$，则数列 $\{a_n\}$ 有界。
\end{theorem}
% UNIT S02.U03
\begin{proof}
由 $a_n\to L$，取 $\varepsilon=1$。则存在正整数 $N$，使得当 $n\geq N$ 时，$|a_n-L|<1$。于是当 $n\geq N$ 时，$|a_n|\leq|a_n-L|+|L|<1+|L|$。令 $M$ 为有限数集 $\{|a_1|,\ldots,|a_{N-1}|,1+|L|\}$ 的最大值；若 $N=1$，则取 $M=1+|L|$。于是对所有 $n\in\mathbb{N}^{+}$ 都有 $|a_n|\leq M$，所以数列有界。
\end{proof}
% UNIT S02.U04
\begin{theorem}
\label{thm:limit-sign}设 $a_n\to L$。若 $L>0$，则存在正整数 $N$，使得当 $n\geq N$ 时 $a_n>0$；若 $L<0$，则存在正整数 $N$，使得当 $n\geq N$ 时 $a_n<0$。特别地，若 $a_n\geq0$ 对所有 $n$ 成立，则 $L\geq0$；若 $a_n\leq0$ 对所有 $n$ 成立，则 $L\leq0$。
\end{theorem}
% UNIT S02.U05
\begin{proof}
当 $L>0$ 时，取 $\varepsilon=L/2$。由极限定义，存在正整数 $N$，使得 $n\geq N$ 时 $|a_n-L|<L/2$。因此 $a_n>L/2>0$。当 $L<0$ 时，取 $\varepsilon=|L|/2$，则当 $n$ 足够大时，$a_n<L+|L|/2=L/2<0$。若 $a_n\geq0$ 且 $L<0$，则这与充分大的 $n$ 满足 $a_n<0$ 矛盾，所以 $L\geq0$。另一种情形同理。
\end{proof}
% UNIT S02.U06
\begin{theorem}
\label{thm:limit-order}设 $a_n\to L$，$b_n\to M$，且存在正整数 $N_0$，使得对所有 $n\geq N_0$ 都有 $a_n\leq b_n$，则 $L\leq M$。
\end{theorem}
% UNIT S02.U07
\begin{proof}
反设 $L>M$，取 $\varepsilon=(L-M)/3>0$。由 $a_n\to L$ 和 $b_n\to M$，存在正整数 $N_1,N_2$，使得当 $n\geq N_1$ 时 $|a_n-L|<\varepsilon$，当 $n\geq N_2$ 时 $|b_n-M|<\varepsilon$。当 $n\geq\max\{N_0,N_1,N_2\}$ 时，$a_n>L-\varepsilon$ 且 $b_n<M+\varepsilon$。所以 $a_n-b_n>L-M-2\varepsilon=(L-M)/3>0$，即 $a_n>b_n$，与题设矛盾。因此 $L\leq M$。
\end{proof}
% UNIT S02.U08
\begin{theorem}
\label{thm:limit-sum-product}设 $a_n\to L$，$b_n\to M$，且 $c\in\mathbb{R}$ 为常数，则 $ca_n\to cL$，$a_n+b_n\to L+M$，以及 $a_nb_n\to LM$。从而 $a_n-b_n\to L-M$。
\end{theorem}
% UNIT S02.U09
\begin{proof}
先证常数倍法则。若 $c=0$，则对任意 $\varepsilon>0$，取 $N=1$，当 $n\geq N$ 时有 $|ca_n-cL|=0<\varepsilon$。若 $c\neq0$，则对任意 $\varepsilon>0$，由 $a_n\to L$ 可取正整数 $N$，使得当 $n\geq N$ 时 $|a_n-L|<\varepsilon/|c|$，于是 $|ca_n-cL|=|c|\,|a_n-L|<\varepsilon$。所以 $ca_n\to cL$。
再证和法则。任给 $\varepsilon>0$，由两个数列的收敛性，分别取正整数 $N_1,N_2$，使得当 $n\geq N_1$ 时 $|a_n-L|<\varepsilon/2$，当 $n\geq N_2$ 时 $|b_n-M|<\varepsilon/2$。取 $N=\max\{N_1,N_2\}$，则当 $n\geq N$ 时，由三角不等式，$|(a_n+b_n)-(L+M)|\leq|a_n-L|+|b_n-M|<\varepsilon/2+\varepsilon/2=\varepsilon$。故 $a_n+b_n\to L+M$。由已证的常数倍法则，$-b_n\to-M$，再用和法则即得 $a_n-b_n\to L-M$。
对乘积，先由 \ref{thm:convergent-bounded} 知道收敛数列有界，取常数 $K>0$，使得对所有 $n$ 都有 $|a_n|\leq K$。利用恒等式 $a_nb_n-LM=a_n(b_n-M)+M(a_n-L)$，有 $|a_nb_n-LM|\leq K|b_n-M|+|M|\,|a_n-L|$。任给 $\varepsilon>0$，由收敛性可分别取正整数 $N_3,N_4$，使得当 $n\geq N_3$ 时 $|b_n-M|<\varepsilon/(2K)$，当 $n\geq N_4$ 时 $|a_n-L|<\varepsilon/[2(|M|+1)]$。取 $N=\max\{N_3,N_4\}$，则当 $n\geq N$ 时，$|a_nb_n-LM|<\varepsilon/2+|M|\varepsilon/[2(|M|+1)]<\varepsilon$。因此由极限定义，$a_nb_n\to LM$。
\end{proof}
% UNIT S02.U10
\begin{theorem}
\label{thm:limit-quotient}设 $a_n\to L$，$b_n\to M$，且 $M\neq0$。若对所有 $n\in\mathbb{N}^{+}$ 都有 $b_n\neq0$，则 $a_n/b_n\to L/M$。
\end{theorem}
% UNIT S02.U11
\begin{proof}
由 $b_n\to M$ 且 $M\neq0$，取 $\varepsilon=|M|/2$。根据极限定义，当 $n$ 足够大时，$|b_n-M|<|M|/2$，从而 $|b_n|\geq|M|-|b_n-M|>|M|/2$。于是 $|1/b_n-1/M|=|b_n-M|/(|M|\,|b_n|)\leq2|b_n-M|/|M|^2\to0$。所以 $1/b_n\to1/M$。再由乘积法则，$a_n/b_n=a_n(1/b_n)\to L(1/M)=L/M$。
\end{proof}
% UNIT S02.U12
\begin{example}
计算下列数列的极限，并说明所使用的运算法则，其中各小问均取 $n\in\mathbb{N}^{+}$：
\[
(1)\quad a_n=3(1/n)+2;\qquad
(2)\quad b_n=\frac{2n+1}{n+3};\qquad
(3)\quad c_n=\left|\frac3n-2\right|.
\]
\end{example}
% UNIT S02.U13
\par\noindent\textbf{解：}(1) 由前一节的基本极限 $1/n\to0$，利用常数倍法则与和法则，
\[
\lim_{n\to\infty}a_n=3\cdot0+2=2.
\]

(2) 对于 $b_n$，由于 $n\in\mathbb{N}^{+}$，可将分子、分母同除以 $n$，得到
\[
b_n=\frac{2+1/n}{1+3/n}.
\]
由 $1/n\to0$，有 $2+1/n\to2$，且 $1+3/n\to1\neq0$。另外，$1+3/n>0$，所以商式有定义。由商法则，
\[
\lim_{n\to\infty}b_n=\frac21=2.
\]
使用商法则前必须确认分母极限非零。

(3) 所用的绝对值极限法则是：若实数数列 $x_n\to L$，其中 $L\in\mathbb{R}$，则 $|x_n|\to|L|$。其依据是反三角不等式
\[
\bigl||x_n|-|L|\bigr|\leq|x_n-L|.
\]
事实上，任给 $\varepsilon>0$，由 $x_n\to L$，存在正整数 $N$，使得当 $n\geq N$ 时 $|x_n-L|<\varepsilon$，于是 $\bigl||x_n|-|L|\bigr|<\varepsilon$，由极限定义即得上述法则。

应用时，先由 $1/n\to0$ 及常数倍、差法则得到 $3/n-2\to-2$，再用绝对值极限法则得到
\[
\lim_{n\to\infty}c_n=\lim_{n\to\infty}|3/n-2|=|-2|=2.
\]
这说明处理含绝对值的数列极限时，可以先求绝对值内部数列的极限，再对所得极限取绝对值。\par
% UNIT S02.U14
\begin{exercise}
\label{mathagent:exercise:8760977b7c41a7ab2d82304ea81f33749e9976ddafc1f87583b542c4b8793254}
设 $a_n\to2$，$b_n\to-1$，且对所有 $n\in\mathbb{N}^{+}$ 都有 $b_n\neq0$。求下列极限，并说明所使用的定理及关键条件：$(1)\ 4a_n-3b_n$；$(2)\ a_nb_n$；$(3)\ (a_n+1)/b_n$。
\end{exercise}
\section{夹逼定理与常见极限}\label{sec:S03}
% UNIT S03.U01
\begin{definition}
若实数数列 $\{a_n\}$ 满足 $a_n\to0$，则称它为无穷小数列。“无穷小”描述整个数列的变化趋势，并不是指某个固定的很小的正数。例如 $1/n$ 是无穷小数列，而常数数列 $a_n=10^{-6}$ 不是无穷小数列。本节将用上下界控制误差，从而处理不能直接使用四则运算法则的极限。
\end{definition}
% UNIT S03.U02
\begin{theorem}
\label{thm:squeeze-sequence}夹逼定理：设 $\{a_n\}$、$\{b_n\}$、$\{c_n\}$ 为实数数列。若存在正整数 $N_0$，使得对所有 $n\geq N_0$ 都有
\[
a_n\leq b_n\leq c_n,
\]
且 $a_n\to L$、$c_n\to L$，其中 $L\in\mathbb{R}$，则 $b_n\to L$。
\end{theorem}
% UNIT S03.U03
\begin{proof}
任给 $\varepsilon>0$。由 $a_n\to L$，存在正整数 $N_1$，使得 $n\geq N_1$ 时 $|a_n-L|<\varepsilon$；由 $c_n\to L$，存在正整数 $N_2$，使得 $n\geq N_2$ 时 $|c_n-L|<\varepsilon$。取
\[
N=\max\{N_0,N_1,N_2\}.
\]
对每个 $n\geq N$，有
\[
L-\varepsilon<a_n\leq b_n\leq c_n<L+\varepsilon.
\]
所以 $|b_n-L|<\varepsilon$。由定义 \ref{def:epsilon-N}，$b_n\to L$。证明同时说明：夹逼关系只需从某个下标开始成立，有限个初始项不影响结论。
\end{proof}
% UNIT S03.U04
夹逼定理常写成误差估计的形式：设 $L\in\mathbb{R}$，$d_n\to0$，并且存在正整数 $N_0$，使得对所有 $n\geq N_0$ 都有
\[
|b_n-L|\leq d_n.
\]
于是这些下标上必有 $d_n\geq0$，且
\[
L-d_n\leq b_n\leq L+d_n.
\]
由和、差法则及定理 \ref{thm:squeeze-sequence}，可得 $b_n\to L$。这不是预先假设 $b_n$ 收敛，而是借助误差上界证明其收敛。若两侧极限不同，则仅凭夹在两者之间不能断言中间数列收敛。\par
% UNIT S03.U05
\begin{theorem}
\label{thm:basic-reciprocal-limit}对于正整数下标 $n$，有
\[
\lim_{n\to\infty}\frac1n=0.
\]
更一般地，对每个固定的正整数 $k$，都有
\[
\lim_{n\to\infty}\frac1{n^k}=0.
\]
\end{theorem}
% UNIT S03.U06
\begin{proof}
任给 $\varepsilon>0$，取正整数 $N>1/\varepsilon$。当 $n\geq N$ 时，$0<1/n\leq1/N<\varepsilon$，故 $1/n\to0$。对任意固定的正整数 $k$，由于 $n\geq1$ 时 $n^k\geq n$，故
\[
0<\frac1{n^k}\leq\frac1n.
\]
由定理 \ref{thm:squeeze-sequence} 得 $1/n^k\to0$。这里 $k$ 是固定的，不能把这一表述不加说明地当作关于任意变化指数的结论。
\end{proof}
% UNIT S03.U07
\begin{theorem}
\label{thm:bounded-times-zero}有界数列与无穷小数列的乘积是无穷小数列。具体地，设 $\{u_n\}$、$\{v_n\}$ 为实数数列，存在常数 $M>0$，使得对每个 $n\in\mathbb{N}^{+}$ 都有 $|u_n|\leq M$，且 $v_n\to0$，则
\[
u_nv_n\to0.
\]
这里不要求 $\{u_n\}$ 本身收敛。
\end{theorem}
% UNIT S03.U08
\begin{proof}
由 $v_n\to0$ 的定义，对任意 $\varepsilon>0$，存在正整数 $N$，使得 $n\geq N$ 时 $|v_n|<\varepsilon/M$。因此对所有 $n\geq N$，
\[
|u_nv_n|=|u_n|\,|v_n|\leq M|v_n|<\varepsilon,
\]
所以 $u_nv_n\to0$。从夹逼的角度看，$|v_n|\to0$ 直接来自极限定义，而
\[
-M|v_n|\leq u_nv_n\leq M|v_n|
\]
的两侧都趋于零。
\end{proof}
% UNIT S03.U09
\begin{example}
设 $n\in\mathbb{N}^{+}$。求数列
\[
a_n=\frac{(-1)^n n}{n^2+1}
\]
的极限，并说明为什么不能把 $(-1)^n$ 和 $n/(n^2+1)$ 的极限直接相乘。
\end{example}
% UNIT S03.U10
\par\noindent\textbf{解：}对每个正整数 $n$，分母 $n^2+1>0$，且
\[
|a_n|=\frac{n}{n^2+1}\leq\frac{n}{n^2}=\frac1n.
\]
由定理 \ref{thm:basic-reciprocal-limit} 及误差估计形式的夹逼定理，$a_n\to0$。

因子 $(-1)^n$ 不收敛。事实上，若它收敛于某个实数 $L$，取 $\varepsilon=1/2$，则充分大的偶数下标与奇数下标分别给出 $|1-L|<1/2$ 和 $|-1-L|<1/2$，从而
\[
2=|1-(-1)|\leq|1-L|+|L+1|<1,
\]
矛盾。因此不能使用要求两个因子都收敛的乘积法则。不过 $|(-1)^n|=1$，且 $n/(n^2+1)\to0$，所以也可用定理 \ref{thm:bounded-times-zero} 得到答案。\par
% UNIT S03.U11
\begin{example}
对 $n\in\mathbb{N}^{+}$，计算
\[
\sqrt{n^2+3}-n
\]
的极限。要求先作代数变形，再用上下界证明，不使用函数极限或导数。
\end{example}
% UNIT S03.U12
\par\noindent\textbf{解：}由于 $n\geq1$，根号内 $n^2+3>0$，且 $\sqrt{n^2+3}+n>0$。乘以共轭式可得
\[
\sqrt{n^2+3}-n
=\frac{(n^2+3)-n^2}{\sqrt{n^2+3}+n}
=\frac3{\sqrt{n^2+3}+n}.
\]
又因 $\sqrt{n^2+3}\geq n$，所以
\[
0<\sqrt{n^2+3}-n\leq\frac3{2n}.
\]
由 $3/(2n)\to0$ 和定理 \ref{thm:squeeze-sequence}，得到
\[
\lim_{n\to\infty}\bigl(\sqrt{n^2+3}-n\bigr)=0.
\]
这里不能把原式中两项分别取实数极限再相减；两项都没有有限实数极限。共轭变形把相减的式子转化成了容易估计的正分式。\par
% UNIT S03.U13
\begin{example}
对 $n\in\mathbb{N}^{+}$，求
\[
r_n=\frac{3n^2-n+2}{2n^2+5}
\]
的极限，并解释分子、分母同除以 $n^2$ 的作用。
\end{example}
% UNIT S03.U14
\par\noindent\textbf{解：}因为 $n>0$，且 $2n^2+5>0$，所以可写成
\[
r_n=\frac{3-1/n+2/n^2}{2+5/n^2}.
\]
定理 \ref{thm:basic-reciprocal-limit} 给出 $1/n\to0$ 和 $1/n^2\to0$。由和、常数倍法则，分子趋于 $3$，分母趋于 $2\neq0$；同时 $2+5/n^2>0$。故由定理 \ref{thm:limit-quotient}，
\[
\lim_{n\to\infty}r_n=\frac32.
\]
同除以 $n^2$，使最高次项变成常数，其余项变成常数乘以 $1/n$ 或 $1/n^2$。这解释了本例为何由最高次项的系数决定极限，而不是把“只看最高次项”当作无需验证条件的规则。\par
% UNIT S03.U15
本节的方法选择可以概括为：遇到有界但不收敛的因子，优先估计绝对值；遇到根式相减，先考虑共轭变形；遇到有理式，先用合适的 $n$ 的幂归一化，再检查商法则的分母条件。夹逼时应写清不等式从哪些下标起成立，并说明两侧具有同一个极限。以下练习的答案统一安排在讲义末尾，本节只列题目。\par
% UNIT S03.U16
\begin{exercise}
\label{mathagent:exercise:719875ae20efa60af241c9d55317b3f3d0138435a92158cbcf8bf58b83fe99e4}
对 $n\in\mathbb{N}^{+}$，用夹逼定理求
\[
b_n=\frac{\sin n}{n+1}
\]
的极限，其中 $\sin n$ 按弧度理解。要求写出对所有正整数 $n$ 成立的绝对值上界，说明上界为何趋于零，并解释为什么不应未经证明就把两个因子的极限相乘。
\end{exercise}
% UNIT S03.U17
\begin{exercise}
\label{mathagent:exercise:2d47f09098678baa1fb51b926189255a253f43fdfbb2aab80fc2d07fb1416ce0}
对 $n\in\mathbb{N}^{+}$，求下列极限，并分别说明所用方法的关键条件：
\[
(1)\quad \sqrt{n^2+n}-n;
\qquad
(2)\quad \frac{2n^2+(-1)^n}{n^2+1}.
\]
第 (1) 题须先作共轭变形；可通过估计 $\sqrt{1+1/n}-1$ 证明 $\sqrt{1+1/n}\to1$，不要直接援引函数连续性。第 (2) 题须明确控制符号变化项产生的误差。
\end{exercise}
\section{单调有界原理与递推数列的极限}\label{sec:S04}
% UNIT S04.U01
\begin{definition}
\label{def:monotone-sequence}设 $\{a_n\}$ 是实数数列。若对所有正整数 $n$ 都有 $a_{n+1}\geq a_n$，则称 $\{a_n\}$ 为单调递增数列；若对所有正整数 $n$ 都有 $a_{n+1}>a_n$，则称其为严格递增数列。若对所有正整数 $n$ 都有 $a_{n+1}\leq a_n$，则称 $\{a_n\}$ 为单调递减数列；若对所有正整数 $n$ 都有 $a_{n+1}<a_n$，则称其为严格递减数列。单调递增或单调递减的数列统称为单调数列。
\end{definition}
% UNIT S04.U02
\begin{definition}
\label{def:bounded-sequence}设 $\{a_n\}$ 是实数数列。若存在实数 $M$，使得对所有正整数 $n$ 都有 $a_n\leq M$，则称 $\{a_n\}$ 有上界；若存在实数 $m$，使得对所有正整数 $n$ 都有 $m\leq a_n$，则称 $\{a_n\}$ 有下界。若同时存在实数 $m,M$，使得对所有正整数 $n$ 都有 $m\leq a_n\leq M$，则称 $\{a_n\}$ 有界。等价地，若存在常数 $C>0$，使得对所有正整数 $n$ 都有 $|a_n|\leq C$，则 $\{a_n\}$ 有界。
\end{definition}
% UNIT S04.U03
\begin{theorem}
\label{thm:monotone-bounded}单调有界原理：若实数数列 $\{a_n\}$ 单调递增且有上界，则 $\{a_n\}$ 收敛；若 $\{a_n\}$ 单调递减且有下界，则 $\{a_n\}$ 收敛。
\end{theorem}
% UNIT S04.U04
\begin{proof}
本证明采用以下已知实数性质作为基础：非空且有上界的实数集合存在上确界（最小上界），非空且有下界的实数集合存在下确界（最大下界）。这是实数完备性的上确界与下确界性质，本节直接使用，不展开其公理体系。先考虑单调递增且有上界的情形。设数列的所有项组成的集合为 $A=\{a_n:n\in\mathbb{N}^{+}\}$。由于 $A$ 非空且有上界，令 $L=\sup A$。任给 $\varepsilon>0$，由上确界的性质，$L-\varepsilon$ 不是 $A$ 的上界，因此存在正整数 $N$，使得 $a_N>L-\varepsilon$。当 $n\geq N$ 时，由数列单调递增，有 $L-\varepsilon<a_N\leq a_n\leq L$。于是 $|a_n-L|<\varepsilon$。根据数列收敛的 $\varepsilon$--$N$ 定义，$a_n\to L$。再考虑单调递减且有下界的情形。令 $B=\{a_n:n\in\mathbb{N}^{+}\}$，并令 $L=\inf B$。任给 $\varepsilon>0$，由下确界的性质，存在正整数 $N$，使得 $a_N<L+\varepsilon$。当 $n\geq N$ 时，由数列单调递减，有 $L\leq a_n\leq a_N<L+\varepsilon$。因此 $|a_n-L|<\varepsilon$，从而 $a_n\to L$。
\end{proof}
% UNIT S04.U05
\begin{theorem}
\label{thm:recursive-limit-method}设 $0<q<1$，$L\in\mathbb{R}$，并由 $a_{n+1}=q a_n+(1-q)L$，$n\in\mathbb{N}^{+}$，定义实数数列 $\{a_n\}$。若 $a_1\leq L$，则 $\{a_n\}$ 单调递增且有上界 $L$，因而收敛于 $L$；若 $a_1\geq L$，则 $\{a_n\}$ 单调递减且有下界 $L$，因而也收敛于 $L$。
\end{theorem}
% UNIT S04.U06
\begin{proof}
先设 $a_1\leq L$。用数学归纳法证明对所有正整数 $n$ 都有 $a_n\leq L$。当 $n=1$ 时结论成立。若 $a_n\leq L$，则由 $0<q<1$ 得 $a_{n+1}=q a_n+(1-q)L\leq qL+(1-q)L=L$。所以 $a_n\leq L$ 对所有正整数 $n$ 成立。接着，由递推关系可得 $a_{n+1}-a_n=q a_n+(1-q)L-a_n=(1-q)(L-a_n)\geq0$。因此数列单调递增，并且由上面的归纳结论知它有上界 $L$。根据定理 \ref{thm:monotone-bounded}，数列收敛。设其极限为 $A$。先说明移位数列也满足 $a_{n+1}\to A$：任给 $\varepsilon>0$，由 $a_n\to A$，存在正整数 $N$，使得对所有 $k\geq N$ 都有 $|a_k-A|<\varepsilon$；当 $n\geq N$ 时，$n+1\geq N$，故 $|a_{n+1}-A|<\varepsilon$。由递推关系和极限的四则运算法则，得到 $A=qA+(1-q)L$。于是 $(1-q)A=(1-q)L$。因为 $1-q\neq0$，所以 $A=L$。当 $a_1\geq L$ 时，同理用归纳法可得 $a_n\geq L$ 对所有正整数 $n$ 成立，且 $a_{n+1}-a_n=(1-q)(L-a_n)\leq0$。因此数列单调递减且有下界 $L$。再次由定理 \ref{thm:monotone-bounded}，数列收敛。设极限为 $A$，由上述移位结论也有 $a_{n+1}\to A$，代入递推关系仍有 $A=qA+(1-q)L$，从而 $A=L$。这里先证明了收敛，之后才可以把极限代入递推关系。
\end{proof}
% UNIT S04.U07
\begin{example}
某商品的价格每期向长期均衡价格调整。设均衡价格为 $40$，初始价格为 $a_1=0$，每期将当前价格与均衡价格之间差距的 $1/2$ 消除，得到递推关系 $a_{n+1}=\frac12a_n+20$，$n\in\mathbb{N}^{+}$。求价格数列的极限，并说明其单调性和有界性。
\end{example}
% UNIT S04.U08
\par\noindent\textbf{解：}将递推关系写成 $a_{n+1}=\frac12a_n+\left(1-\frac12\right)40$。这里 $q=1/2$，满足 $0<q<1$，且 $a_1=0\leq40$。由定理 \ref{thm:recursive-limit-method}，数列 $\{a_n\}$ 单调递增并且有上界 $40$，因此收敛，且 $\lim_{n\to\infty}a_n=40$。也可以直接检查前几步：$a_2=20$，$a_3=30$，$a_4=35$，价格逐期接近 $40$。这些数值只能帮助理解趋势，不能代替对所有正整数 $n$ 的归纳证明。经济含义是：在该调整规则下，价格不会超过均衡价格，并且会逐步接近均衡价格。\par
% UNIT S04.U09
\begin{exercise}
\label{mathagent:exercise:3095db2dc0f723686b84b8afdd22af197f8af13ce063f1be09f10d01d7c714d4}
设数列 $\{a_n\}$ 由 $a_1=10$，$a_{n+1}=\frac12a_n+15$，$n\in\mathbb{N}^{+}$，确定。证明该数列单调递增且有上界，进而求其极限。要求明确写出归纳步骤、单调性证明、有界性证明以及求极限时所列出的方程。
\end{exercise}
\section*{练习答案}
% UNIT S01.U10
\par\noindent\textbf{\hyperref[mathagent:exercise:41706b9e371cbf8283a4f1192e660b4426644ffb8f650af52f47cef5491b1ca0]{练习 \ref*{mathagent:exercise:41706b9e371cbf8283a4f1192e660b4426644ffb8f650af52f47cef5491b1ca0} 的解答：}}给定任意 $\varepsilon>0$，注意到
\[
\left|\frac{2n+1}{n}-2\right|=\left|\frac{1}{n}\right|=\frac{1}{n}.
\]
取正整数 $N>\frac{1}{\varepsilon}$。当 $n\geq N$ 时，$n>0$，从而
\[
\left|\frac{2n+1}{n}-2\right|=\frac{1}{n}\leq\frac{1}{N}<\varepsilon.
\]
因此根据 $\varepsilon$--$N$ 定义，
\[
\lim_{n\to\infty}\frac{2n+1}{n}=2.
\]\par
% UNIT S02.U15
\par\noindent\textbf{\hyperref[mathagent:exercise:8760977b7c41a7ab2d82304ea81f33749e9976ddafc1f87583b542c4b8793254]{练习 \ref*{mathagent:exercise:8760977b7c41a7ab2d82304ea81f33749e9976ddafc1f87583b542c4b8793254} 的解答：}}由常数倍法则与差法则，$4a_n-3b_n\to4\cdot2-3\cdot(-1)=11$。由乘积法则，$a_nb_n\to2\cdot(-1)=-2$。又由和法则，$a_n+1\to3$。由于 $b_n\to-1\neq0$，且题设保证 $b_n\neq0$，可使用商法则，得到 $(a_n+1)/b_n\to3/(-1)=-3$。关键条件是分母数列的极限非零，并且商式在所讨论的下标上有定义。\par
% UNIT S03.U16.answer
\par\noindent\textbf{\hyperref[mathagent:exercise:719875ae20efa60af241c9d55317b3f3d0138435a92158cbcf8bf58b83fe99e4]{练习 \ref*{mathagent:exercise:719875ae20efa60af241c9d55317b3f3d0138435a92158cbcf8bf58b83fe99e4} 的解答：}}对所有正整数 $n$，由 $|\sin n|\leq1$ 及 $n+1>0$，有
\[
|b_n|=\frac{|\sin n|}{n+1}\leq\frac1{n+1}\leq\frac1n.
\]
由于 $1/n\to0$，故
\[
-\frac1n\leq b_n\leq\frac1n
\]
两侧均趋于零。由夹逼定理，$\displaystyle\lim_{n\to\infty}b_n=0$；同时上述估计也说明上界 $1/(n+1)$ 趋于零。

虽然 $b_n=(\sin n)\cdot1/(n+1)$，但乘积极限法则要求两个因子数列均收敛，不能仅凭 $|\sin n|\leq1$ 就认定 $\{\sin n\}$ 收敛，因而不能未经证明就把两个因子的极限相乘。本题只需利用正弦因子的有界性，无需确定其收敛性；也可用有界数列与无穷小数列的乘积趋于零来说明。\par
% UNIT S03.U17.answer
\par\noindent\textbf{\hyperref[mathagent:exercise:2d47f09098678baa1fb51b926189255a253f43fdfbb2aab80fc2d07fb1416ce0]{练习 \ref*{mathagent:exercise:2d47f09098678baa1fb51b926189255a253f43fdfbb2aab80fc2d07fb1416ce0} 的解答：}}(1) 对所有正整数 $n$，根号内为正，且共轭式 $\sqrt{n^2+n}+n>0$，所以
\[
\sqrt{n^2+n}-n
=\frac{n}{\sqrt{n^2+n}+n}
=\frac1{\sqrt{1+1/n}+1}.
\]
这里利用了 $n>0$，从而 $\sqrt{n^2+n}=n\sqrt{1+1/n}$。再作共轭变形并估计：
\[
0\leq\sqrt{1+1/n}-1
=\frac{1/n}{\sqrt{1+1/n}+1}
\leq\frac1{2n}.
\]
由 $1/(2n)\to0$ 和夹逼定理，$\sqrt{1+1/n}-1\to0$，故 $\sqrt{1+1/n}\to1$。变形后分母处处为正，且趋于 $2\neq0$，由商法则得
\[
\lim_{n\to\infty}\bigl(\sqrt{n^2+n}-n\bigr)=\frac12.
\]
关键是先用共轭变形及夹逼证明根式数列的极限，再检查商法则的分母条件，并未援引函数连续性。

(2) 由于 $n^2+1>0$，将原式与 $2$ 作差得
\[
\frac{2n^2+(-1)^n}{n^2+1}-2
=\frac{(-1)^n-2}{n^2+1}.
\]
符号变化项满足 $|(-1)^n|=1$，因此对所有正整数 $n$，
\[
\left|\frac{2n^2+(-1)^n}{n^2+1}-2\right|
\leq\frac{|(-1)^n|+2}{n^2+1}
=\frac3{n^2+1}\leq\frac3{n^2}\to0.
\]
由误差估计形式的夹逼定理，
\[
\lim_{n\to\infty}\frac{2n^2+(-1)^n}{n^2+1}=2.
\]
关键是分母为正，且符号变化项造成的误差具有趋于零的绝对值上界；无需对 $(-1)^n$ 本身取极限。\par
% UNIT S04.U10
\par\noindent\textbf{\hyperref[mathagent:exercise:3095db2dc0f723686b84b8afdd22af197f8af13ce063f1be09f10d01d7c714d4]{练习 \ref*{mathagent:exercise:3095db2dc0f723686b84b8afdd22af197f8af13ce063f1be09f10d01d7c714d4} 的解答：}}先证明对所有正整数 $n$ 都有 $a_n\leq30$。当 $n=1$ 时，$a_1=10\leq30$。若 $a_n\leq30$，则 $a_{n+1}=\frac12a_n+15\leq\frac12\cdot30+15=30$。由数学归纳法，得到 $a_n\leq30$，$n\in\mathbb{N}^{+}$。所以数列有上界 $30$。再由递推关系计算相邻两项之差：$a_{n+1}-a_n=\frac12a_n+15-a_n=\frac12(30-a_n)\geq0$。因此 $\{a_n\}$ 单调递增。数列单调递增且有上界，根据定理 \ref{thm:monotone-bounded}，它收敛。设 $\lim_{n\to\infty}a_n=A$。由于已经证明数列收敛，才可以在递推关系中取极限，得到 $A=\frac12A+15$。所以 $\frac12A=15$，$A=30$。因此 $\lim_{n\to\infty}a_n=30$。\par
\end{document}
